\section{假设检验、单调似然比与预警功效曲线}
\label{sec:part11}

\subsection{暴发预警的统计检验构造与单调似然比}

在传染病监测预警中（估计目标 1），决策者最核心的关切不是参数的渐近均方误差，而是及时判断疫情是否处于扩张阶段：
\begin{equation}
  H_0\colon R_t \le 1 \quad \text{versus} \quad H_1\colon R_t > 1.
  \label{eq:hypo-def}
\end{equation}

\begin{theorembox}{定理 11.1（负二项分布族的严格单调似然比性质）}
在假设 H1--H4 下，给定祖代规模 $n$，观测病例数 $C_t$ 的条件分布族关于参数 $R_t$ 具有严格单调似然比（Monotone Likelihood Ratio, MLR）。
\end{theorembox}

\begin{proof}
任取两个参数值 $R_2 > R_1 > 0$。根据式~\eqref{eq:der-ct-dist}，观测值 $C_t = c$ 在两个参数下的概率质量函数之比为：
\begin{align}
  \Lambda(c; R_2, R_1) &\equiv \frac{P(C_t = c \mid R_2)}{P(C_t = c \mid R_1)} \nonumber \\
  &= \frac{\frac{\Gamma(c+nk)}{c!\,\Gamma(nk)} \left(\frac{k}{k+\rho R_2}\right)^{nk} \left(\frac{\rho R_2}{k+\rho R_2}\right)^c}{\frac{\Gamma(c+nk)}{c!\,\Gamma(nk)} \left(\frac{k}{k+\rho R_1}\right)^{nk} \left(\frac{\rho R_1}{k+\rho R_1}\right)^c} \nonumber \\
  &= \left( \frac{k+\rho R_1}{k+\rho R_2} \right)^{nk} \left[ \frac{\rho R_2 / (k+\rho R_2)}{\rho R_1 / (k+\rho R_1)} \right]^c \nonumber \\
  &= \left( \frac{k+\rho R_1}{k+\rho R_2} \right)^{nk} \left[ \frac{R_2(k+\rho R_1)}{R_1(k+\rho R_2)} \right]^c.
  \label{eq:lr-ratio}
\end{align}
考察以 $c$ 为指数的底数项：
\begin{equation}
  \Theta(R_2, R_1) \equiv \frac{R_2(k+\rho R_1)}{R_1(k+\rho R_2)} = \frac{k R_2 + \rho R_1 R_2}{k R_1 + \rho R_1 R_2}.
\end{equation}
由于 $R_2 > R_1$ 且 $k > 0$，比较分子与分母：
\begin{equation}
  (k R_2 + \rho R_1 R_2) - (k R_1 + \rho R_1 R_2) = k(R_2 - R_1) > 0.
\end{equation}
因此分子严格大于分母，底数满足：
\begin{equation}
  \Theta(R_2, R_1) > 1.
\end{equation}
由于底数大于 1，比值函数 $\Lambda(c; R_2, R_1)$ 是统计量 $c$ 的\textbf{严格单调递增函数}。
证毕。
\end{proof}

\subsection{Karlin--Rubin 定理与一致最优（UMP）精确检验}

根据数理统计中的 Karlin--Rubin 定理 \citep{karlin1956theory, lehmann2005testing}，对于具有单调似然比性质的单参数分布族，单侧假设检验的一致最优势检验（Uniformly Most Powerful, UMP）必为临界域形式。

对于离散型负二项分布，由于概率累积存在阶梯跳跃，标准 Karlin--Rubin UMP 检验在理论上包含边界随机化：当 $C_t > c$ 时拒绝，当 $C_t = c$ 时以概率 $\gamma \in [0, 1)$ 随机化拒绝，从而使得实际水平严格等于名义水平 $\alpha$。在公共卫生监测工程实践中，为了避免人为抛硬币随机化导致的判断不确定性，通常采用确定性的保守非随机化检验规则。

\begin{derivationbox}{推导 11.1：临界值与精确功效函数}
设名义显著性水平为 $\alpha \in (0, 1)$。定义临界值 $c_\alpha$ 为使得零假设临界点 $R_t = 1$ 下第一类错误概率不超过 $\alpha$ 的最小整数：
\begin{equation}
  c_\alpha \equiv \min \left\{ c \in \mathbb{N} \ \Big|\ P_{R_t = 1}(C_t \ge c) \le \alpha \right\}.
\end{equation}
本文采用非随机化确定性规则：当且仅当 $C_t \ge c_\alpha$ 时拒绝原假设 $H_0$。
\begin{enumerate}
  \item \textbf{实际水平（Size）：}由于功效函数在单调似然比族下关于 $R_t$ 单调非减，原假设区域内的最大犯错概率在边界 $R_t = 1$ 处取得：
  \begin{equation}
    \alpha^* \equiv \sup_{R_t \le 1} P_{R_t}(C_t \ge c_\alpha) = P_{R_t = 1}(C_t \ge c_\alpha) \le \alpha.
  \end{equation}
  非随机化检验 $\{C_t \ge c_\alpha\}$ 对应随机化检验在 $c = c_\alpha - 1, \gamma = 1$ 的特例，因此它在其实际达到的显著性水平 $\alpha^*$ 下是一致最优（UMP）的；若在边界 $c_\alpha - 1$ 处允许随机化，则可使实际水平严格等于名义水平 $\alpha$ 并取得更高功效。
  \item \textbf{精确功效（Power）：}当真实传播发生暴发，即备择假设点 $R_t = 1 + \delta$（$\delta > 0$）时，观测病例数服从负二项分布：
  \begin{equation}
    C_t \sim \NB\!\left(nk,\ \frac{k}{k + \rho(1+\delta)}\right).
  \end{equation}
  其精确检测功效定义为拒绝原假设的概率，由负二项分布累积分布函数 $F_{\NB}$ 直接给出：
  \begin{equation}
    \beta(n, \delta) \equiv P_{R_t = 1+\delta}(C_t \ge c_\alpha) = 1 - F_{\NB}\left( c_\alpha - 1;\ nk,\ \frac{k}{k + \rho(1+\delta)} \right).
    \label{eq:exact-power-proof}
  \end{equation}
\end{enumerate}
\end{derivationbox}

\subsection{检验功效锯齿震荡与 Wald 检验的尺寸膨胀}

在实际工程应用中，许多流行病学软件常使用大样本正态近似的 Wald 检验统计量：
\begin{equation}
  Z \equiv \frac{\widehat{R}_t - 1}{\sqrt{V_0}} \sim \mathcal{N}(0, 1), \qquad V_0 \equiv \Var(\widehat{R}_t \mid R_t = 1) = \frac{1}{\rho n} + \frac{1}{kn}.
\end{equation}

\begin{insightbox}{数值病态洞见 11.1：Wald 检验严重虚假报警机制}
由于负二项分布在样本量较小及存在超传播（小 $k$）时呈现极端的右偏重尾特征，正态近似在尾部严重失真：
\begin{itemize}
  \item 在基准参数 $k = 0.1, m = 5$ 例的小样本下，名义水平为 5\% 的 Wald 检验，其实际第一类错误概率高达 \textbf{7.81\%}；
  \item 在 $m < 40$ 的区间内，不同超传播强度 $k = 0.1, 0.35, 1.0$ 下 Wald 检验的峰值假阳性率分别达到 \textbf{8.03\%}、\textbf{8.75\%} 和 \textbf{10.40\%}，超出安全标准一倍以上；
  \item 相比之下，基于式~\eqref{eq:exact-power-proof} 的精确检验在全参数网格上的最大实际检验水平为 \textbf{0.049997}，严格被 5\% 上界控制。
\end{itemize}

此外，由于临界值 $c_\alpha$ 随样本规模 $n$ 呈现阶梯状整数跳跃，精确功效曲线 $\beta(n, \delta)$ 关于 $n$ 表现出剧烈的\textbf{锯齿状震荡（Serration）}：
例如，在 $k = 0.35, \delta = 0.3$ 时，当祖代规模 $n$ 从 400 增加至 699 的过程中，由于每当临界点 $c_\alpha$ 增加 1，拒绝阈值突然上抬，导致检测功效在此过程中\textbf{逆向下降了 82 次}！

因此，监测系统必须区分两个规模标准：
\begin{enumerate}
  \item 首次达到 80\% 功效的脆弱规模；
  \item 越过最后一个锯齿跌落、此后稳定保持在 80\% 功效之上的鲁棒规模。
\end{enumerate}
在 $\rho = 0.25, k = 0.35$ 下，检出 15\% 暴发增长（$\delta = 0.15$）所需的两种祖代规模分别为 2{,}071 人与 2{,}107 人，对应预期病例数 $m \approx 527$ 例，约为方差平衡规模 $\mx = 53.1$ 的整整十倍。
\end{insightbox}
